\subsection{Sensory selection in DendriNet populations}
\label{sec:dendrinet_selection}

The independent seven-compartment implementation isolates the gate mechanism.
To test its transfer to the production network implementation, we trained
DendriNet populations with fixed, grouped input contacts. The principal
configuration contained sixteen somata, each with a $[4,2]$ tree: eight
terminal compartments, four proximal compartments and one soma. Each pair
of terminals received two features of one sensory stream. An affine readout
combined the sixteen somatic voltages. Leak and reversal potentials were
those of the artificial conductance models in Methods. All excitatory,
inhibitory and coupling conductances, and the readout, were plastic; input
supports were fixed rather than chosen by adaptive contact selection.

\paragraph{Tasks and forward controls.}
For four streams indexed by $b=0,1,2,3$, independent latent features
$z_{bj}\sim\mathcal U[-2,2]$, $j=1,2$, supplied positive terminal inputs
$\exp(z_{bj})$ and $\exp(-z_{bj})$. A one-hot cue $\bm c$ selected one stream.
Every terminal also received the same four excitatory cue channels, so none
of the forward controls lacked context information. The targets were
specified independently of the network:
\begin{equation}
y=\sum_{b=0}^3 c_b(-1)^b
\left[\frac{\tanh z_{b1}+\tanh z_{b2}}2
       +\beta\tanh z_{b1}\tanh z_{b2}\right].
\label{eq:si_selection_target}
\end{equation}
Alternating signs prescribe opposed feature tuning across streams. They do
not supply targets to the network. The primary task sets $\beta=0$ and uses
identity compartment outputs. A secondary condition sets $\beta=0.25$ and
uses $\tanh(V_p)$ at proximal compartments, allowing the two features to
interact. This secondary condition changes both target and representation;
it is not an isolated test of the effect of adding a nonlinearity.

In the shunting forward model, proximal inhibitory activity was
$a_b^{\rm I}=4(1-c_b)$. With $D_p^0=1+\sum_i g_{i\to p}^{\rm den}$ and child current
$I_p=\sum_i g_{i\to p}^{\rm den}a_i$, the proximal voltage was
$V_p^{\rm sh}=I_p/(D_p^0+g_p^{\rm I}a_b^{\rm I})$. Tonic inhibition replaced
$a_b^{\rm I}$ by its training-distribution mean $\bar a^{\rm I}=3$; the excitatory cue inputs
were unchanged. A current-injection control preserved that tonic denominator
but varied a context-dependent current:
\begin{equation}
V_p^{\rm cur}=
\frac{I_p-g_p^{\rm I}(a_b^{\rm I}-\bar a^{\rm I})\bar V_p}
     {D_p^0+g_p^{\rm I}\bar a^{\rm I}}.
\end{equation}
Here $\bar V_p$ was the mean tonic voltage at initialization, estimated
using training inputs and then held fixed. This matches the incremental
inhibitory current at that reference voltage, not the full voltage
distribution throughout training. Shunting, tonic and current models had
the same parameter count, cue channels and paired initial parameters.

\paragraph{Learning rules.}
The local rules detached intercompartment connections when differentiating
the loss, while retaining each compartment's exact local derivatives.
The affine readout supplied a separate somatic error to each neuron:
\begin{equation}
\widehat y=\sum_{u=1}^{16}w_u^{\rm out}V_u+b^{\rm out},\qquad
\ell=\frac{(\widehat y-y)^2}{2\operatorname{Var}_{\rm train}(y)},\qquad
\delta_u^V=\frac{w_u^{\rm out}(\widehat y-y)}{\operatorname{Var}_{\rm train}(y)}.
\end{equation}
Training averaged $\ell$ over the minibatch. Thus broadcast here means
sharing a neuron's error across its own compartments, not sharing one
scalar across all sixteen neurons.
Proximal and somatic updates used this error without spatial attenuation.
At a terminal whose proximal parent was $p$, the relative-resistance rule
instead delivered
\begin{equation}
\widehat\delta_{ui}^V=\delta_u^V h_{up},\qquad
h_{up}=\frac{D_{up}^0}{D_{up}^0+g_{up}^{\rm I}a_b^{\rm I}}.
\label{eq:si_dendrinet_gate}
\end{equation}
Unit broadcast set $h=1$. A wrong-branch control cyclically permuted
the four parent gates within each neuron. A uniform-gain control replaced
them by their within-neuron root-mean-square,
$h_u^{\rm rms}=(\sum_{p=1}^4h_{up}^2/4)^{1/2}$, preserving
the gate-vector norm but not the norm of the eligibility-weighted parameter
update. Exact BP provided a reference. All five rules shared the same
forward model within a condition and trained the same parameters.
In the tonic and current controls, the local gate still used the supplied
trial-dependent inhibitory cue. These conditions therefore test a proposed
credit signal without its corresponding forward shunt, not the exact
resistance of those alternative forward models. In the nonlinear-parent
condition the gate omitted the parent's activation derivative. Local
training accessed neither exact path products nor exact-gradient diagnostics.

The chain rule explains this distinction. In the shunting model, exact
terminal credit is
\begin{equation}
\delta_{ui}^V=\delta_u^V\,
\frac{g_{p\to u}^{\rm den}g_{i\to p}^{\rm den}}{D_uD_p^0}
h_{up} f'_p(V_p),\qquad
D_u=1+\sum_p g_{p\to u}^{\rm den}.
\end{equation}
At fixed weights, the prefactor is constant across examples. With identity
parent outputs, $f'_p=1$ and the relative-resistance gate preserves the
entire context dependence of terminal transport, but not its branch-specific
constant factors. With nonlinear parents, $f'_p=1-\tanh^2(V_p)$ varies with
both input and context. Resistance alone no longer supplies all of the
example-dependent transport. This predicts a limitation of the particular
gate, not an inability of local signals in general to support nonlinear
dendritic learning.

\paragraph{Distractor stress test and development.}
The held-out stress test replaced only irrelevant latent inputs by
$z_{bj}^{(s)}=[1+(s-1)(1-c_b)]z_{bj}$, with
$s\in\{1,1.5,2,2.5,3\}$. Relevant features, cue and target were identical
across severities within a seed. Because the input encoding is exponential,
this manipulation broadens the presynaptic-rate distribution nonlinearly.
The ordinary test set was an independent draw from the training distribution;
it was not the $s=1$ member of the stress-test sample. This is a synthetic
robustness test, not a natural-image benchmark.

Three development seeds (2026092100--2026092102) compared Adam rates
0.01, 0.03 and 0.1 for each forward/rule combination over 2,048 updates,
giving 180 trajectories. Minimum mean validation NMSE selected one rate
per condition. The fresh cohort used twenty disjoint paired seeds
(2026092200--2026092219), the selected rates and 4,096 updates: 400
trajectories over the three identity-output forward models and the
shunting nonlinear-parent condition. Training, validation and test sizes
were 2,048, 1,024 and 4,096; batches contained 128 examples. Validation
every 128 updates, including initialization, selected the checkpoint.
NMSE divided by the variance of targets in the evaluation split.
Log conductances were clipped to $[-9,9]$ and gradient norms to ten;
bound and clipping contacts were recorded. Rates and checkpoints were
not selected using test or distractor-stress outcomes.

Two primary contrasts were specified before fresh outcomes: broadcast
minus relative-resistance NMSE, and uniform-RMS minus relative-resistance
NMSE, both at $s=3$ in the identity-output shunting condition. Two-sided
paired sign-flip tests were Holm-adjusted over these two contrasts;
95\% percentile intervals resampled whole seeds. Other comparisons were
descriptive. Common-state diagnostics measured the core-gradient cosine
and the change in NMSE after an equal-length, $10^{-3}$ displacement in
log-conductance space, holding the readout fixed. These diagnostics neither
generated nor selected the local training updates. The protocol followed
the exploratory cohort below and was internally frozen, not externally
preregistered.

\paragraph{Fresh-seed outcomes and boundary.}
Table~\ref{tab:dendrinet_selection} reports all fresh conditions. In the
separable shunting task at severity three, the paired broadcast-minus-gate
NMSE difference was 0.01038 [0.00864, 0.01244], and the uniform-RMS-minus-gate
difference was 0.00959 [0.00816, 0.01120]. Both were positive in all twenty
seeds; both Holm-adjusted sign-flip $P$ values were $3.81\times10^{-6}$.
Mean errors were 0.010506 for broadcast, 0.009716 for uniform RMS,
0.0001213 for the relative-resistance gate and 0.0001191 for exact BP.
The near agreement of the last two means is not an equivalence test.
Wrong-branch gating gave 0.06017. Ordinary-test errors were small for all
three principal rules, so the most pronounced benefit concerns robustness
outside the training input range, rather than rescue of ordinary task learning.

The forward controls separated this credit benefit from input suppression.
At severity three, exact BP gave mean NMSE 0.2435 with tonic inhibition and
0.3056 with matched current injection, compared with 0.0001191 with
forward shunting. Relative-resistance gating gave 0.2420 and 0.2250 in the
two alternative forwards. All twenty paired differences favored shunting
for each of the three principal rules. These comparisons concern trained
forward models; they do not hold their final weights fixed. Nor do the
differences in differences imply a synergistic interaction: the absolute
broadcast-minus-gate reduction was larger in the poorly performing
tonic and current controls.

At common exact-trained shunting states, the mean core-gradient cosine with
exact BP was 0.312 for broadcast, 0.484 for relative-resistance gating and
0.416 for uniform RMS gating. Mean NMSE decreases after equal-length core
steps were $1.02\times10^{-6}$, $1.57\times10^{-6}$ and
$1.37\times10^{-6}$, respectively, with the readout fixed. These descriptive
diagnostics support a directional contribution beyond a change in overall
update magnitude; they do not establish equivalence of training trajectories.

A secondary common-rate check was specified during fresh execution,
without altering the primary protocol. Sixty additional trajectories on
the same twenty seeds supplied the three missing rule--rate combinations,
allowing broadcast, uniform RMS and the spatial gate to be compared at both
Adam 0.03 and 0.1. The gate reduced distractor error in all twenty seeds
against each control at each rate. The broadcast-minus-gate differences
were 0.01031 [0.00857, 0.01236] at 0.03 and
0.000681 [0.000543, 0.000849] at 0.1. Thus the effect persisted at common
rates but its magnitude was rate-sensitive; these repeated seed blocks
are not a second independent confirmation.

The nonlinear-parent interaction condition gave a different result.
Exact BP reached mean ordinary-test NMSE $5.06\times10^{-5}$, whereas
relative-resistance gating gave 0.05109, broadcast 0.04712 and uniform
RMS gating 0.04893. This task is learnable by the nonlinear network under
exact credit, but the tested resistance-only approximation did not recover
that performance. The contrast is consistent with its missing parent-gain
factor; the experiment does not isolate that factor from every other
optimization difference. Bound contacts also limit interpretation: in the
separable shunting cohort, 12/20 exact and 7/20 gated trajectories contacted
a bound, compared with none for broadcast or uniform RMS. In the nonlinear
condition the respective counts were 2/20, 18/20, 19/20 and 19/20. These
counts cover full trajectories, not only selected checkpoints. Tonic and
current controls also frequently contacted bounds. No claim of unbounded
or asymptotically optimal learning follows.

All 400 primary-cohort checkpoints reproduced their ordinary-test and five
stress-test NMSEs exactly on replay. All sixty common-rate checkpoints
reproduced those outcomes to absolute tolerance $10^{-10}$. Complete
endpoints, validation trajectories, paired contrasts, diagnostics and source
hashes are retained in
\path{source_data/curated_publication/inhibitory_selection_*}; generators
and verification code are in \path{scripts/inhibitory_selection/}.

\paragraph{Exploratory cohort and its representation limit.}
An earlier matrix crossed one or sixteen somata, two or four streams
using configurations $(1,2)$, $(16,2)$ and $(16,4)$, selection or graded
mixture targets, three forward models, five rules and two rates
(0.01 and 0.03). Three paired development seeds gave 540 trajectories
of 1,024 updates with log-conductance bounds $[-7,7]$. Mixture cues
were sampled from $\operatorname{Dirichlet}(1,\ldots,1)$ and weighted
all streams in Eq.~\ref{eq:si_selection_target}. All compartments had
identity outputs, but the target included $\beta=0.25$.

That configuration cannot represent its within-stream product term:
conditional on the cue, each terminal depends on one latent feature,
proximal denominators depend only on context, and the soma and readout
combine feature functions additively. For independent symmetric inputs,
the product is orthogonal to every such additive predictor. Writing
$v=\mathbb E[\tanh^2z]=1-\tanh(2)/2$, the population squared-error floor
normalized by target variance is
\begin{equation}
\frac{\beta^2v^2}{v/2+\beta^2v^2}
=\frac{v}{8+v}=0.06081\quad(\beta=0.25).
\end{equation}
The same ratio applies to the independent-stream mixture after averaging
the squared cue weights. It is not label noise or evidence of failed credit
assignment. The primary follow-up removes this term; the nonlinear-parent
condition instead removes the additive-model obstruction, without assuming
exact representability of the entire target.

The pilot motivated the robustness question. In the sixteen-neuron,
four-stream shunting condition, mean selection NMSE at $s=3$ was 0.0655
with relative-resistance gating, 0.0830 with broadcast, 0.0826 with uniform
RMS gating and 0.1475 with wrong-branch gating; exact BP gave 0.0653.
The same local gate did not improve graded mixtures: its NMSE was 0.1669,
versus 0.1622 with broadcast and 0.1571 with uniform RMS gating.
These are descriptive means over three development seeds, not confirmation
or evidence of equivalence. Of 90 selected rates, 89 were the grid maximum;
92 of 270 selected runs chose the training cap and 50 contacted a bound.
The full matrix, including the unfavorable outcomes, is retained in
\path{source_data/curated_publication/inhibitory_transfer_pilot_*}.
